\chapter{K4510 BASIC}\label{cha:kbasic}

K4510 BASIC is the machine's BASIC for writing programs. It speaks
EhBASIC's words (PRINT, INPUT, FOR, GOSUB, PLOT ...), but it is
\emph{compiled}: \verb|k4510-bas| turns your \verb|.BAS| file into Mad Pascal, and Mad
Pascal turns that into a \verb|.prg|, the same kind of program as everything
in \verb|/APPS|. The result runs many times faster than EhBASIC.

You write a program in PROG (or VI, or any editor), save it as
\verb|NAME.BAS|, and then:

\begin{center}
\small
\begin{tabular}{@{}l >{\raggedright\arraybackslash}p{0.72\linewidth}@{}}
\toprule
\textbf{Where} & \textbf{What to do} \\
\midrule
At the prompt & \texttt{BAS NAME} compiles NAME.BAS into \verb|name.prg|; \verb|NAME| runs it \\
In PROG & \textbf{F9} compiles; \textbf{Ctrl-F9} compiles and runs \\
\bottomrule
\end{tabular}
\end{center}

There is no ``immediate mode'' (typing \texttt{PRINT 2+2} and getting 4). For a
quick calculation or a one-line try, use RX: there \verb|PRINT| works like
\verb|SAY|, so \texttt{PRINT 6*7} prints 42.

EhBASIC is still on the machine for the old numbered programs in
\verb|/LANG/EHBASIC|. Every one of its examples is in \verb|/LANG/BASIC/EX| too,
written again in K4510 BASIC: \verb|DEMOS| and \verb|BENCH| are the menus.

\section{A first program}

\begin{type}
' guess the number -- the first K4510 BASIC program
CONST TRIES = 7
secret = 42
n% = 0
DO
  INPUT "Your guess"; g
  n% = n% + 1
  IF g < secret THEN
    PRINT "Higher"
  ELSEIF g > secret THEN
    PRINT "Lower"
  END IF
LOOP UNTIL g = secret OR n% = TRIES
GOSUB Bravo
PRINT "Tries:"; n%; "of"; TRIES
END
Bravo:
  IF g = secret THEN PRINT "You got it!" ELSE PRINT "Too bad"
  RETURN
\end{type}

No line numbers, keywords in any case, a label (\verb|Bravo:|) where an old
BASIC would have had \texttt{GOSUB 1000}.

\section{Writing it down}

\begin{itemize}
\item \textbf{One statement per line}, or several separated by \verb|:|.
\item \textbf{Keywords in any case}: \verb|print|, \verb|Print| and \verb|PRINT| are the same. So are variable names: \verb|Score| and \verb|SCORE| are one variable.
\item \textbf{Comments} start with \verb|'| or \verb|REM| and run to the end of the line.
\item \textbf{Labels}: a name followed by \verb|:| at the start of a line (\verb|Again:|). \texttt{GOTO Again} and \texttt{GOSUB Again} jump there.
\item \textbf{Line numbers} still work, as labels: \texttt{100 PRINT "HI"} and \texttt{GOTO 100}. You can mix both, and leave most lines without one.
\end{itemize}

\section{Variables}

\begin{center}
\small
\begin{tabular}{@{}l >{\raggedright\arraybackslash}p{0.56\linewidth} l@{}}
\toprule
\textbf{Written} & \textbf{Holds} & \textbf{Starts as} \\
\midrule
\verb|score| & a number with decimals (single precision, about 7 figures) & 0 \\
\verb|lives%| & a whole number, -32768 to 32767 & 0 \\
\verb|name$| & text, up to 255 characters & ``'' \\
\bottomrule
\end{tabular}
\end{center}

\verb|a|, \verb|a%| and \verb|a$| are three different variables. Names can be as long
as you like, and every letter counts (EhBASIC only looked at the first
two).

A whole-number variable rounds what it is given: \texttt{n\% = 3.7} makes 4. A
value outside -32768..32767 stops the program with ``Overflow''.

\textbf{Every variable belongs to the whole program}, wherever it is used,
except the parameters of a SUB or FUNCTION and what a SUB or FUNCTION
DIMs for itself (see below).

\subsection{Arrays}

\begin{type}
DIM scores(10), board%(7, 7), names$(30)
scores(3) = 9.5
board%(0, 7) = 1
\end{type}

\begin{itemize}
\item One or two dimensions. \texttt{DIM a(10)} has 11 elements, 0 to 10.
\item The size is a whole number or a CONST.
\item An array used without a DIM is \verb|DIM|med to 10 for you.
\item An element of a text array holds up to 80 characters.
\item An array may use about 24 KB at most (DIM tells you if it is bigger).
\item A wrong index (\verb|scores(11)|, \verb|scores(-1)|) stops the program with ``Index out of range''.
\end{itemize}

\subsection{CONST}

\begin{type}
CONST TRIES = 7
CONST GREETING$ = "Hello"
\end{type}

A CONST is a variable that may not be changed: \texttt{TRIES = 8} is an error
when the program is compiled. A CONST made from a plain number can be
an array's size: \texttt{DIM grid(SIZE)}.

\section{Values and operators}

\begin{center}
\small
\begin{tabular}{@{}>{\raggedright\arraybackslash}p{0.44\linewidth} >{\raggedright\arraybackslash}p{0.44\linewidth}@{}}
\toprule
\textbf{Operator} & \textbf{Means} \\
\midrule
\texttt{+ -{} * /} & add, subtract, multiply, divide (\verb|/| always gives decimals) \\
\verb|\| & whole-number division: \texttt{7 \textbackslash{} 2} is 3 \\
\verb|MOD| & remainder: \texttt{7 MOD 2} is 1 \\
\verb|^| & power: \texttt{2 \textasciicircum{} 10} is 1024 \\
\verb|+| on text & joins: \texttt{"ab" + "cd"} is \verb|"abcd"| \\
\texttt{= <> < > <= >=} & compare numbers, or text (alphabetical, by character code) \\
\texttt{AND OR XOR NOT} & combine conditions; on numbers, they work bit by bit \\
\bottomrule
\end{tabular}
\end{center}

\verb|^| comes first, then a minus sign, then \texttt{* /}, \verb|\|, \verb|MOD|, \texttt{+ -{}}, the
comparisons, \verb|NOT|, \verb|AND|, \verb|XOR|, \verb|OR|. So \texttt{-{}2 \textasciicircum{} 2} is -4. Brackets
change the order as usual.

A comparison used as a number is -1 (true) or 0 (false):
\texttt{x = (3 > 2)} makes x -1. A number used as a condition is true when it
is not 0: \texttt{IF lives\% THEN ...}.

Numbers can be written \verb|12|, \verb|3.75|, \verb|.5|, \verb|1.5E+12| or in hexadecimal
\verb|&HFF|.

\section{Statements}

\subsection{Printing and asking}

\begin{center}
\small
\begin{tabular}{@{}>{\raggedright\arraybackslash}p{0.44\linewidth} >{\raggedright\arraybackslash}p{0.44\linewidth}@{}}
\toprule
\textbf{Statement} & \textbf{What it does} \\
\midrule
\texttt{PRINT a; b\$; c} & prints; \verb|;| keeps going on the same line \\
\texttt{PRINT a, b} & \verb|,| moves to the next column of 14 \\
\texttt{PRINT TAB(20); "x"} & \verb|TAB(n)| moves to column n, \verb|SPC(n)| prints n spaces \\
\texttt{PRINT "no new line";} & a \verb|;| or \verb|,| at the end keeps the cursor on the line \\
\texttt{? "hi"} & \verb|?| is short for PRINT \\
\texttt{INPUT "Your age"; age} & prints the question and \texttt{? }, waits for a line \\
\texttt{INPUT "Name: ", n\$} & with \verb|,| no \texttt{? } is added \\
\texttt{INPUT a, b} & several values, typed separated by commas \\
\texttt{GET k\$} & the key pressed, or ``'' if none: does not wait \\
\verb|CLS| & clears the screen \\
\texttt{COLOR 7} / \texttt{COLOR 7, 6} & text colour, and background (0-15, the palette) \\
\texttt{LOCATE 10, 30} & moves the cursor to row 10, column 30 \\
\bottomrule
\end{tabular}
\end{center}

A number is printed with a space in front (where a minus sign would
go) and one after: \texttt{PRINT 5; 6} shows \texttt{ 5  6 }. Text is printed as it
is. Numbers show at most six figures, without useless zeros: \verb|3.75|,
\verb|0.333333|, and very big or very small ones as \verb|1.5E+12|.

\subsection{Deciding}

\begin{type}
IF x > 10 THEN PRINT "big" ELSE PRINT "small"

IF age < 13 THEN
  PRINT "child"
ELSEIF age < 18 THEN
  PRINT "teenager"
ELSE
  PRINT "adult"
END IF

SELECT CASE n%
  CASE 1, 2
    PRINT "one or two"
  CASE 3 TO 9
    PRINT "a few"
  CASE IS >= 10
    PRINT "lots"
  CASE ELSE
    PRINT "none"
END SELECT
\end{type}

\begin{itemize}
\item The one-line IF keeps everything after THEN on that line; \texttt{IF x THEN 100} and \texttt{IF x GOTO 100} jump to a label or line number.
\item A many-line IF has nothing after THEN and ends with \texttt{END IF}.
\item \texttt{SELECT CASE} works with numbers and with text (\texttt{CASE "yes", "y"}).
\end{itemize}

\subsection{Repeating}

\begin{type}
FOR i = 1 TO 10 STEP 2 ... NEXT i
WHILE lives% > 0 ... WEND
DO WHILE x < 5 ... LOOP
DO UNTIL done% ... LOOP
DO ... LOOP UNTIL key$ = "q"
DO ... LOOP WHILE x < 5
DO ... LOOP              ' forever, until EXIT DO
\end{type}

\begin{itemize}
\item \verb|STEP| may be negative (\texttt{FOR i = 10 TO 1 STEP -{}1}) or have decimals. Without STEP, the count goes up by 1.
\item \verb|NEXT| may name its variable (\texttt{NEXT i}) or not.
\item \texttt{EXIT FOR} and \texttt{EXIT DO} leave the loop they are in.
\end{itemize}

\subsection{Jumping}

\begin{center}
\small
\begin{tabular}{@{}>{\raggedright\arraybackslash}p{0.44\linewidth} >{\raggedright\arraybackslash}p{0.44\linewidth}@{}}
\toprule
\textbf{Statement} & \textbf{What it does} \\
\midrule
\texttt{GOTO label} & continues at the label (or line number) \\
\texttt{GOSUB label} & runs from the label until RETURN, then comes back \\
\verb|RETURN| & back to the statement after the GOSUB \\
\verb|END| or \verb|STOP| & ends the program \\
\bottomrule
\end{tabular}
\end{center}

GOSUBs can be inside each other up to 64 deep.

\subsection{SUB and FUNCTION}

\begin{type}
SUB Stars (n%)
  FOR i% = 1 TO n%: PRINT "*";: NEXT
  PRINT
END SUB

FUNCTION Area (w, h)
  Area = w * h
END FUNCTION

Stars 5
CALL Stars(3)
PRINT Area(3, 4)
\end{type}

\begin{itemize}
\item A SUB is called by its name (with or without brackets round the values) or with \verb|CALL|.
\item A FUNCTION gives back the value assigned to its own name.
\item They can be written anywhere in the file, before or after they are used.
\item \texttt{EXIT SUB} and \texttt{EXIT FUNCTION} leave early.
\item \textbf{Local variables:} \texttt{DIM x, n\%, s\$} or \texttt{DIM t(3)} inside a SUB or FUNCTION makes them its own, 0 or ``'' at every call, hiding the program's variables of the same name. Its parameters are its own too. Every other variable inside it is the program's variable of that name.
\item A SUB or FUNCTION cannot call itself (no recursion): its own variables are kept in one place, not one set per call.
\item \texttt{Name: next statement} is a call of the SUB Name followed by another statement; for any other name, \verb|Name:| at the start of a line is a label.
\item Labels, GOTO and GOSUB stay outside SUBs and FUNCTIONs.
\end{itemize}

\subsection{DATA and READ}

\begin{type}
FOR i = 1 TO 3: READ name$, age: PRINT name$; age: NEXT
DATA Ada, 36, "Grace", 85
DATA Alan, 41
RESTORE           ' READ starts again from the first DATA
RESTORE Colours   ' ... or from the DATA after a label
\end{type}

\begin{itemize}
\item DATA can be anywhere in the program; READ takes the items in order.
\item An item is a number, a ``quoted'' text, or text written as it is (the spaces round it trimmed).
\item READ past the last item stops the program with ``Out of DATA''.
\end{itemize}

\subsection{ON ... GOTO, DEF FN}

\begin{type}
ON choice% GOTO One, Two, Three     ' 1 to One, 2 to Two ...
ON choice% GOSUB One, Two, Three
DEF FNarea(w, h) = w * h            ' a one-line function
DEF FNfull$(a$, b$) = a$ + " " + b$
PRINT FNarea(3, 4)
\end{type}

A value outside 1 to the number of labels goes on to the next statement.

\subsection{The shell: * and SHELL}

\begin{type}
*MODE 640x480               ' the rest of the line goes to K/OS, : and all
SHELL "MODE " + size$       ' a command worked out by the program
SHELL "INVADERS"            ' runs another program; this one goes on after it
\end{type}

A program that changes the screen gets the shell's own screen back when
it ends (which clears it): wait for a key before END if the last words
matter.

\subsection{Other statements}

\begin{center}
\small
\begin{tabular}{@{}>{\raggedright\arraybackslash}p{0.44\linewidth} >{\raggedright\arraybackslash}p{0.44\linewidth}@{}}
\toprule
\textbf{Statement} & \textbf{What it does} \\
\midrule
\texttt{LET x = 5} & the same as \texttt{x = 5} \\
\texttt{SLEEP 1.5} & waits that many seconds \\
\verb|RANDOMIZE| & new random numbers (every program already starts with fresh ones) \\
\texttt{POKE address, value} & writes a byte into memory: 0-65535 is what the CPU sees, I/O included; 65536 and up is the machine's far memory (to 16 MB), as \verb|PEEK| reads it \\
\bottomrule
\end{tabular}
\end{center}

\section{Functions}

\begin{center}
\small
\begin{tabular}{@{}>{\raggedright\arraybackslash}p{0.46\linewidth} >{\raggedright\arraybackslash}p{0.46\linewidth}@{}}
\toprule
\textbf{Function} & \textbf{Gives} \\
\midrule
\verb|ABS(x)| & x without its sign \\
\verb|INT(x)| & x rounded down: \verb|INT(-2.5)| is -3 \\
\verb|FIX(x)| & x with its decimals cut off: \verb|FIX(-2.5)| is -2 \\
\verb|SGN(x)| & -1, 0 or 1 \\
\verb|SQR(x)| & square root \\
\texttt{SIN COS TAN ATN (x)} & trigonometry, in radians \\
\verb|EXP(x)| \verb|LOG(x)| & e to the power x; natural logarithm \\
\verb|RND| or \verb|RND(1)| & a random number from 0 up to (not including) 1 \\
\verb|LEN(a$)| & how many characters \\
\verb|ASC(a$)| & the code of the first character (0 if empty) \\
\verb|CHR$(n)| & the character with code n \\
\verb|VAL(a$)| & the number written in a\$ (\verb|VAL("3.5")| is 3.5) \\
\verb|STR$(x)| & the number as text, without the leading space \\
\texttt{LEFT\$(a\$, n)} \texttt{RIGHT\$(a\$, n)} & the first / last n characters \\
\texttt{MID\$(a\$, start)} \texttt{MID\$(a\$, start, n)} & n characters from position start (from 1) \\
\texttt{INSTR(a\$, b\$)} & where b\$ is in a\$ (1 = first character), 0 if it is not \\
\verb|UCASE$(a$)| \verb|LCASE$(a$)| & in capitals / small letters \\
\verb|SPACE$(n)| & n spaces \\
\verb|INKEY$| & the key pressed, or ``'': does not wait \\
\verb|PEEK(address)| & a byte of memory \\
\bottomrule
\end{tabular}
\end{center}

To make a whole number from 1 to 6: \texttt{INT(RND * 6) + 1}.

\section{Graphics}

The graphics words draw on a picture laid over the text, as EhBASIC's
do. Colour 0 is see-through: the text shows through it.

\begin{center}
\small
\begin{tabular}{@{}>{\raggedright\arraybackslash}p{0.46\linewidth} >{\raggedright\arraybackslash}p{0.46\linewidth}@{}}
\toprule
\textbf{Statement} & \textbf{What it does} \\
\midrule
\texttt{GRAPHICS 1} & the picture on, 320 x 240 (MODE 2), as in EhBASIC \\
\texttt{GRAPHICS 2} & the picture on, 640 x 480 (MODE 0) \\
\texttt{GRAPHICS 3} & the picture on, at the size of the screen as it is (after \texttt{*MODE 400x300}, say) \\
\texttt{GRAPHICS 0} & the picture off, and the text mode GRAPHICS found back \\
\texttt{PALETTE i, r, g, b} & colour i (0-255) as red, green, blue (0-255); 0-15 are the text's \\
\texttt{SPRDEF n, page, w, h, bpp} & sprite n (0-127) has its shape at page*256 in far memory, w x h (8, 16, 32, 64), 4 or 8 bits a pixel \\
\texttt{SPRITE n, x, y} & puts sprite n there and shows it \\
\texttt{SPROFF n} & hides it \\
\verb|GCLS| & clears the picture \\
\texttt{PLOT x, y, c} & one dot \\
\texttt{LINE x1, y1, x2, y2, c} & a line \\
\texttt{BOX x1, y1, x2, y2, c} & a filled rectangle \\
\texttt{TRI x1, y1, x2, y2, x3, y3, c} & a filled triangle \\
\texttt{CIRCLE x, y, r, c} & a circle \\
\verb|GWIDTH| \verb|GHEIGHT| & the picture's width and height, in dots \\
\bottomrule
\end{tabular}
\end{center}

The first drawing word turns the picture on (\texttt{GRAPHICS 3}) if the
program has not. A program that ends with the picture on waits for a
key, then takes it away. Coordinates and colours are whole numbers;
(0, 0) is the top left.

\section{Errors}

\textbf{When compiling.} A mistake stops the compile with a message and the
line it is on: ``IF (END IF) is never closed'', ``there is no label
Bravo'', ``TRIES is a CONST: it cannot change'', ``a number was expected
here, not a string''. In PROG, F9 puts the cursor on that line.

\textbf{When running.} Some mistakes can only be seen while the program
runs. The program stops, prints the line and the reason, and goes back
to K/OS; PROG then puts the cursor on that line:

\begin{center}
\small
\begin{tabular}{@{}>{\raggedright\arraybackslash}p{0.46\linewidth} >{\raggedright\arraybackslash}p{0.46\linewidth}@{}}
\toprule
\textbf{Message} & \textbf{Because} \\
\midrule
\texttt{Line N: Division by zero} & \verb|/|, \verb|\| or \verb|MOD| by 0 \\
\texttt{Line N: Index out of range (DIM it bigger?)} & an array index outside its DIM \\
\texttt{Line N: Overflow: a whole number (\%) holds -{}32768 to 32767} & too big for a \verb|%| variable \\
\texttt{Line N: RETURN without GOSUB} & a RETURN reached without a GOSUB \\
\texttt{Line N: Too many GOSUBs inside each other} & more than 64 GOSUBs not yet RETURNed \\
\bottomrule
\end{tabular}
\end{center}

\section{Compared with EhBASIC and QuickBASIC}

\textbf{Like EhBASIC:} the words, the graphics statements, the variable
types, PRINT's look.

\textbf{Unlike EhBASIC:} compiled, not interpreted; no line numbers needed;
keywords in any case; every letter of a name counts; multi-line IF,
SELECT CASE, SUB and FUNCTION, labels, LOCATE, COLOR, CLS, INKEY\$.

\textbf{Like QuickBASIC:} the block statements (IF ... END IF, DO ... LOOP,
SELECT CASE), labels, SUB and FUNCTION, CONST.

\textbf{Not there (yet):}

\begin{itemize}
\item \texttt{PRINT USING}, \texttt{LINE INPUT}, \verb|WRITE|
\item \verb|SWAP|, \verb|ERASE|, \verb|REDIM|
\item recursion; \verb|SHARED|, \verb|STATIC|
\item labels, \verb|GOTO| and \verb|GOSUB| inside a SUB or FUNCTION
\item files (\verb|OPEN|, \verb|LOAD|, \verb|SAVE|), \verb|SOUND|, \verb|TIMER|, \verb|DATE$|
\item arrays of more than two dimensions; text array elements longer than 80
\item immediate mode (use RX)
\end{itemize}
